% ═══════════════════════════════════════════════════════════════
% MT-05a-ML: MUSTERLÖSUNG MINITEST LOS HOBBYS
% Spanisch Klasse 7 - Sequenz 5a
% ═══════════════════════════════════════════════════════════════
% Punkte: 26 BE
% ═══════════════════════════════════════════════════════════════

\documentclass[11pt, a4paper]{article}

\newif\ifswmodus
\swmodusfalse

\usepackage[utf8]{inputenc}
\usepackage[T1]{fontenc}
\usepackage[ngerman]{babel}
\usepackage{helvet}
\renewcommand{\familydefault}{\sfdefault}

\usepackage[margin=2cm]{geometry}
\usepackage{enumitem}
\usepackage{tabularx}
\usepackage{booktabs}

\usepackage{xcolor}
\usepackage{tcolorbox}
\tcbuselibrary{skins}

\usepackage{fancyhdr}
\usepackage{lastpage}
\usepackage{needspace}

% FARBEN
\definecolor{spanisch-color}{HTML}{c41e3a}
\definecolor{grau-color}{HTML}{888888}
\definecolor{hellgrau-color}{HTML}{f5f5f5}
\definecolor{loesung-color}{HTML}{c62828}

\definecolor{sw-dark}{HTML}{000000}
\definecolor{sw-gray}{HTML}{666666}
\definecolor{sw-white}{HTML}{ffffff}

\ifswmodus
    \colorlet{spanisch}{sw-dark}
    \colorlet{grau}{sw-gray}
    \colorlet{hellgrau}{sw-white}
    \colorlet{loesung}{sw-dark}
\else
    \colorlet{spanisch}{spanisch-color}
    \colorlet{grau}{grau-color}
    \colorlet{hellgrau}{hellgrau-color}
    \colorlet{loesung}{loesung-color}
\fi

\newcommand{\fachfarbe}{spanisch}

\newcommand{\thema}{Minitest: Los hobbys}
\newcommand{\fach}{Spanisch}
\newcommand{\klasse}{7}
\newcommand{\maxpunkte}{26}
\newcommand{\zeit}{15 Minuten}
\newcommand{\datum}{\today}

\pagestyle{fancy}
\fancyhf{}
\renewcommand{\headrulewidth}{0.4pt}
\renewcommand{\footrulewidth}{0.4pt}

\lhead{\fach{} -- Klasse \klasse}
\chead{\textcolor{loesung}{\textbf{MT-05a MUSTERLÖSUNG}}}
\rhead{\datum}

\lfoot{\zeit}
\cfoot{}
\rfoot{Seite \thepage\ von \pageref{LastPage}}

\newcommand{\punktebox}[1]{\hfill\fbox{\small \_/#1}}
\newcommand{\luecke}[1]{\rule{#1}{0.4pt}}
\newcommand{\lsg}[1]{\textcolor{loesung}{\textbf{#1}}}

\newtcolorbox{infobox}{
    colback=hellgrau,
    colframe=\fachfarbe,
    boxrule=1pt,
    arc=0pt,
    left=8pt, right=8pt, top=8pt, bottom=8pt
}

\newenvironment{aufgabe}[1]{%
    \needspace{5\baselineskip}%
    \nopagebreak[4]%
    \vspace{10pt}
    \noindent\textbf{\textcolor{\fachfarbe}{Aufgabe #1}}
    \nopagebreak[4]%
    \vspace{5pt}
    \nopagebreak[4]%
    \begin{list}{}{\leftmargin=0pt}
    \item[]
}{%
    \end{list}
}

\newcommand{\es}[1]{\textcolor{\fachfarbe}{\textbf{#1}}}

\begin{document}

% ═══════════════════════════════════════════════════════════════
% KOPFBEREICH
% ═══════════════════════════════════════════════════════════════

\begin{center}
    {\Large\bfseries\textcolor{\fachfarbe}{\thema}}\\[0.3em]
    {\large\textcolor{loesung}{\textbf{--- MUSTERLÖSUNG ---}}}
\end{center}

\vspace{5pt}

\noindent
\begin{tabularx}{\textwidth}{|X|X|X|}
\hline
\textbf{Name:} \lsg{Musterschüler/in} &
\textbf{Klasse:} \klasse &
\textbf{Datum:} \datum \\
\hline
\end{tabularx}

\vspace{10pt}

\begin{infobox}
\textbf{Zeit:} \zeit{} | \textbf{Punkte:} \maxpunkte{} BE | \textbf{Hilfsmittel:} keine
\end{infobox}

\vspace{15pt}

% ═══════════════════════════════════════════════════════════════
% AUFGABE 1: Zuordnung (6P)
% ═══════════════════════════════════════════════════════════════

\begin{aufgabe}{1 -- Zuordnung}
Verbinde die spanischen Hobbys mit der deutschen Bedeutung. \punktebox{6}
\end{aufgabe}

\begin{center}
\begin{tabular}{rcl}
\es{jugar al fútbol} & \lsg{---} & schwimmen \lsg{(4)} \\[0.8em]
\es{escuchar música} & \lsg{---} & Musik hören \lsg{(2)} \\[0.8em]
\es{ver la tele} & \lsg{---} & Fahrrad fahren \lsg{(5)} \\[0.8em]
\es{nadar} & \lsg{---} & fernsehen \lsg{(3)} \\[0.8em]
\es{montar en bici} & \lsg{---} & Fußball spielen \lsg{(1)} \\[0.8em]
\es{leer} & \lsg{---} & lesen \lsg{(6)} \\
\end{tabular}
\end{center}

\textit{\textcolor{grau}{Je 1P pro korrekter Zuordnung. Lösung: 1-Fußball spielen, 2-Musik hören, 3-fernsehen, 4-schwimmen, 5-Fahrrad fahren, 6-lesen}}

\vspace{15pt}

% ═══════════════════════════════════════════════════════════════
% AUFGABE 2: Lückentext (8P)
% ═══════════════════════════════════════════════════════════════

\begin{aufgabe}{2 -- Lückentext}
Ergänze die Sätze mit dem passenden Hobby auf Spanisch. \punktebox{8}
\end{aufgabe}

\noindent
a) Ich spiele gerne Gitarre: Me gusta \lsg{tocar la guitarra}\\[1em]
b) Ich chatte gerne mit Freunden: Me gusta \lsg{chatear}\\[1em]
c) Ich spiele gerne Videospiele: Me gusta \lsg{jugar a videojuegos}\\[1em]
d) Ich schwimme gerne: Me gusta \lsg{nadar}

\textit{\textcolor{grau}{Je 2P pro korrekter Antwort. Orthografie beachten!}}

\vspace{15pt}

% ═══════════════════════════════════════════════════════════════
% AUFGABE 3: Frage und Antwort (6P)
% ═══════════════════════════════════════════════════════════════

\begin{aufgabe}{3 -- Frage und Antwort}
Beantworte die Fragen auf Spanisch. \punktebox{6}
\end{aufgabe}

\noindent
a) Wie fragst du auf Spanisch: ,,Was machst du gerne?''

\lsg{¿Qué te gusta hacer?}

\vspace{1em}

\noindent
b) Wie antwortest du: ,,Ich höre gerne Musik''?

\lsg{Me gusta escuchar música.}

\vspace{1em}

\noindent
c) Wie antwortest du: ,,Ich sehe gerne fern''?

\lsg{Me gusta ver la tele.}

\textit{\textcolor{grau}{Je 2P. Auch ,,Me gusta ver la televisión'' akzeptiert.}}

\vspace{15pt}

% ═══════════════════════════════════════════════════════════════
% AUFGABE 4: Übersetzung (6P)
% ═══════════════════════════════════════════════════════════════

\begin{aufgabe}{4 -- Übersetzung}
Übersetze ins Spanische. \punktebox{6}
\end{aufgabe}

\noindent
a) Fußball spielen = \lsg{jugar al fútbol}\\[1em]
b) Fahrrad fahren = \lsg{montar en bici / montar en bicicleta}\\[1em]
c) fernsehen = \lsg{ver la tele / ver la televisión}

\textit{\textcolor{grau}{Je 2P. Artikel beachten: jugar AL fútbol (nicht jugar fútbol).}}

% ═══════════════════════════════════════════════════════════════
% PUNKTEÜBERSICHT
% ═══════════════════════════════════════════════════════════════

\vspace{25pt}
\hrule
\vspace{10pt}

\noindent
\begin{tabularx}{\textwidth}{|X|c|c|c|c||c|}
\hline
\textbf{Aufgabe} & 1 & 2 & 3 & 4 & \textbf{Gesamt} \\
\hline
\textbf{max. Punkte} & 6 & 8 & 6 & 6 & \textbf{26} \\
\hline
\textbf{erreicht} & \lsg{6} & \lsg{8} & \lsg{6} & \lsg{6} & \lsg{26} \\[0.8em]
\hline
\end{tabularx}

\vspace{10pt}

\begin{flushright}
\textbf{Note:} \lsg{14 NP (100\%)}
\end{flushright}

\end{document}
